% !TeX root = ../../Book3.tex
\chapter{Laurent 级数、奇点与留数}
\label{b3:ch:03}
\BookChapterNumber{b3:ch:03}
\begin{chapterabstract}
Laurent 展开把环域中的全纯函数分成正幂与主部，并由主部分类孤立奇点。留数提取积分中的唯一局部贡献，辐角原理和 Rouché 定理则把围道积分变成稳定的零点计数。本章同时区分函数在无穷远的奇点与微分在无穷远的留数。
\end{chapterabstract}

\begin{learninggoals}
\begin{itemize}
\item\label{b3:item:03-goal-1}能在不同环域选择 Laurent 展开并求出系数。
\item\label{b3:item:03-goal-2}能用局部有界性、主部和倒数函数辨认三类孤立奇点。
\item\label{b3:item:03-goal-3}能计算有限点与无穷远留数，并证明实积分附加弧的极限。
\item\label{b3:item:03-goal-4}能以重数使用辐角原理与 Rouché 定理，并解释零点稳定性。
\end{itemize}
\end{learninggoals}

函数 \(1/(z(z-1))\) 在 \(0<|z|<1\) 与 \(|z|>1\) 有不同的展开。负幂并不是错误项，而是中心孔洞留下的信息；其中 \(-1\) 次系数又恰好决定绕中心的积分。这种局部展开与全局计数的关系，是本章从级数走向留数的原因。

\ProseParagraphBreak

本章需要\chapref{b3:ch:02}的积分公式与估计，以及\chapref{b3:ch:01}的孤立零点分解。John B. Conway 的《Functions of One Complex Variable I》\cite{Conway1978Functions}和 Elias M. Stein、Rami Shakarchi 的《Complex Analysis》\cite{Stein2003Complex}涵盖 Laurent 与留数主题，可用于延伸阅读；下面的奇点判别与计数推导在正文展开。

% !TeX root = ../../Book3.tex
\section{Laurent 级数}
\label{b3:sec:0301}

Taylor 级数要求中心点也在定义域内。去掉中心后，负幂可能出现，但它们并不破坏环域内部的一致收敛。把内外两条边界上的 Cauchy 核分别展开，就能同时得到正幂与负幂；两个部分有不同的几何来源。

\subsection{环域}
\label{b3:subsec:030101}

设 \(0\le r<R\le\infty\)，记
\[
 A(a;r,R)\coloneqq\{\,z\in\mathbb C\mid r<|z-a|<R\,\}.
\]
它称为以 \(a\) 为中心的\term[huanyu]{环域}[annulus]；\(r=0\) 时是去心圆盘。所谓 Laurent 级数的收敛，指正幂部分与负幂部分分别收敛，不能用正负指标混合的任意排列改变和。

\ProseParagraphBreak

负幂级数 \(\sum_{n\ge1}c_{-n}(z-a)^{-n}\) 在倒数变量 \(w=1/(z-a)\) 中是没有常数项的幂级数。若该幂级数半径为 \(S\)，它在 \(|z-a|>1/S\) 上收敛；正幂部分则在 \(|z-a|<R\) 上收敛。两者的公共范围因而天然是环域，可能为空。边界圆上的收敛须分别考察。

\ProseParagraphBreak

例如 \(1/(z(z-1))\) 在 \(0<|z|<1\) 时写成 \(-\sum_{n\ge0}z^{n-1}\)，在 \(|z|>1\) 时写成 \(\sum_{n\ge0}z^{-n-2}\)。展开中心相同，而穿过 \(|z|=1\) 时表达式改变，原因是圆周上的奇点 \(z=1\) 阻断了统一展开。

\ProseParagraphBreak

选择展开环域时，可以先将所有与中心的距离画在数轴上。例如函数 \(1/((z-1)(z-2))\) 的奇点到中心0的距离为1与2，因此有三个候选范围：\(|z|<1\)、\(1<|z|<2\) 和 \(|z|>2\)。中间范围中的 \(1/(z-1)\) 应展开成负幂，而 \(1/(z-2)\) 应展开成正幂。同一个有理式里既有正幂又有负幂，并不表示使用了互相矛盾的几何级数；它们的公共收敛域正好是指定环域。

\ProseParagraphBreak

Laurent 主部是否有限，也取决于中心在哪里。函数 \(1/(z-1)\) 在中心0的小圆盘是 Taylor 级数，在中心1则是只有一个负幂的 Laurent 级数，在无穷远坐标又具有可去奇点。分类和系数总要附一个中心，不能只说“这个函数有负幂”而略去展开范围。

\ProseParagraphBreak

若两个奇点到中心距离相同，例如 \(1/(z^2+1)\)，它们在同一分界圆上。尽管圆周上只有两个点坏掉，也不能把圆内与圆外的 Laurent 展开拼成一个穿过该圆的同心环域展开，因为一个 Laurent 级数在其开放收敛环域上处处全纯，会把这两个真极点也包括进去。

\subsection{Laurent 展开}
\label{b3:subsec:030102}

\term[laurent-jishu]{Laurent 级数}[Laurent series]允许整数幂指标。
\begin{theorem}[Laurent]\label{b3:thm:03-laurent}
若 \(f\) 在 \(A(a;r,R)\) 上全纯，则有唯一展开
\[
 f(z)=\sum_{n=-\infty}^{\infty}c_n(z-a)^n,
\]
其中正幂部分在 \(|z-a|<R\) 全纯，负幂部分在 \(|z-a|>r\) 全纯；级数在每个闭子环域上一致绝对收敛。
\end{theorem}
\begin{proof}
固定 \(r<\rho_1<|z-a|<\rho_2<R\)。沿环域的一条径向线切开，将内外圆的边界用两条相反方向的切口边连接；切口可以选择不经过 \(z\)。对切开区域使用 Cauchy 公式，再让两条切口边靠拢。核与函数在切口附近连续，两个切口积分相消，得到
\[
 f(z)=\frac1{2\pi\mathrm{i}}\int_{|\zeta-a|=\rho_2}\frac{f(\zeta)}{\zeta-z}\,\mathrm{d}\zeta-
 \frac1{2\pi\mathrm{i}}\int_{|\zeta-a|=\rho_1}\frac{f(\zeta)}{\zeta-z}\,\mathrm{d}\zeta,
\]
两圆均取逆时针方向。在外圆上，展开 \((\zeta-z)^{-1}=\sum_{n\ge0}(z-a)^n/(\zeta-a)^{n+1}\)；在内圆上，展开 \((\zeta-z)^{-1}=-\sum_{n\ge0}(\zeta-a)^n/(z-a)^{n+1}\)。若 \(z\) 限于严格夹在两圆之间的闭子环域，两个几何级数都一致绝对收敛，故可以逐项积分。

所得系数为 \(c_n=(2\pi\mathrm{i})^{-1}\int_{|\zeta-a|=\rho}f(\zeta)/(\zeta-a)^{n+1}\,\mathrm{d}\zeta\)。在两个容许半径之间，被积函数全纯；对其环形边界切开并用 Cauchy 定理，两个圆周积分相等。因而系数不依赖所选半径。取半径分别靠近 \(r\) 和 \(R\)，前述表示覆盖整个环域。对任意闭子环域，选更内、更外两圆，几何估计给出一致绝对收敛。

若另有一组 Laurent 系数，在某个中间圆周上逐项积分乘以 \((z-a)^{-k-1}\)，只有指标 \(k\) 的项留下 \(2\pi\mathrm{i}\)。于是两组系数逐个相等。
\end{proof}
这里切口不是对函数选分支，而是把有孔区域分解成可以使用 Cauchy 公式的区域。切口两侧的函数本来相同，故会相消；对多值表达式不能直接照用这一步。

\subsection{Laurent 系数}
\label{b3:subsec:030103}

\begin{proposition}[Laurent 系数公式]\label{b3:prop:03-laurent-coefficients}
在 Laurent 定理的条件下，对任意 \(r<\rho<R\) 与整数 \(n\)，
\[
 c_n=\frac1{2\pi\mathrm{i}}\int_{|z-a|=\rho}\frac{f(z)}{(z-a)^{n+1}}\,\mathrm{d}z,
 \qquad |c_n|\le M_\rho\rho^{-n},
\]
其中 \(M_\rho=\max_{|z-a|=\rho}|f(z)|\)。
\end{proposition}
\begin{proof}
公式已由前一定理的构造与唯一性得到。对积分使用长度估计，圆长 \(2\pi\rho\) 与分母 \(\rho^{n+1}\) 相消，即得对正负指标同时有效的估计。
\end{proof}
当 \(n=-m<0\) 时，估计成为 \(|c_{-m}|\le M_\rho\rho^m\)。若函数靠近中心有界，令 \(\rho\downarrow0\) 就迫使全部负幂系数为零。正幂系数由外侧增长控制，负幂系数由内侧增长控制；同一公式的两个方向将给出奇点分类。

\ProseParagraphBreak

级数仍可在闭子环域上逐项微分，因为两部分分别是在原变量与倒数变量中的幂级数；微分后 \((z-a)^n\) 产生 \(n(z-a)^{n-1}\)，包括负指标。逐项原函数则需额外检查 \(n=-1\) 项，它不能在整个环域上取单值原函数。

\ProseParagraphBreak

Laurent 展开中的系数可由对称性提前筛选。若 \(f(-z)=f(z)\) 且定义环域以0为中心，唯一性给 \(c_n=(-1)^n c_n\)，所以奇数指标系数为零，包括负的奇数指标。若 \(f(\overline z)=\overline{f(z)}\)，则把展开共轭并替换变量后，由唯一性得到全部系数为实数。这里的对称必须在整个环域成立；只在实轴知道函数值实，并且有内部聚点时，也可由恒等定理传播，但要先保证所比较的共轭反射函数全纯。

\ProseParagraphBreak

负幂项不能用普通原函数遗漏。若 \(f\) 在同心环域的留数为0，可以对全部其他项逐项积分，得到在该环域收敛的 Laurent 原函数。其正幂与负幂部分仍在闭子环域一致收敛，可用幂级数估计核对。若留数非零，绕中心圆周积分非零，端点公式立即排除原函数。因此对同心环域上的全纯函数，留数为零恰是存在单值原函数的条件；对于有多个孔洞的区域，要分别检查各孔的周期。

\begin{proposition}\label{b3:prop:03-annulus-primitive}
设 \(f\) 在环域 \(r<|z-a|<R\) 全纯，Laurent 系数为 \(c_n\)。它在该环域有单值全纯原函数，当且仅当 \(c_{-1}=0\)。
\end{proposition}
\begin{proof}
有原函数时任意中间圆周的积分为零，所以 \(2\pi\mathrm{i}c_{-1}=0\)。反之，去掉 \(n=-1\) 项后令
\[
 F(z)=\sum_{n\ne-1}\frac{c_n}{n+1}(z-a)^{n+1}.
\]
在任意闭子环域上，正幂与负幂部分均正规收敛；略扩大子环域后用几何估计可同时控制导数级数。因此逐项求导合法，且 \(F'=f\)。
\end{proof}
这里每个负幂项 \((z-a)^{-n}\)、\(n\ge2\) 都有原函数；只有 \((z-a)^{-1}\) 会产生对数的多值性。故“奇点存在”并不等于“没有原函数”。\(1/z^2\) 虽在0有极点，却在去心平面有原函数 \(-1/z\)。闭路积分测量的是留数这一特定系数，不能直接读出全部主部。

% !TeX root = ../../Book3.tex
\section{孤立奇点}
\label{b3:sec:0302}

奇点分类要先确认中心确实被一个完整的去心圆盘隔离。设 \(a\in\mathbb C\)、\(R>0\)，若 \(f\) 在 \(0<|z-a|<R\) 上全纯，而暂未在 \(a\) 定义，就称 \(a\) 为 \(f\) 的\term[guli-qidian]{孤立奇点}[isolated singularity]。若其他奇点不断积聚到 \(a\)，便不存在这样的去心圆盘，不能将它们全部当作中心的一个孤立奇点。

\ProseParagraphBreak

Laurent 展开把问题变成负幂系数的判别。没有负幂时，正幂级数会填入中心；只有有限个非零负幂时，乘上适当的中心距离幂就能消去分母；负幂无限多时，这两种修补都不再有效。下面分别把这三种情形转成函数值及局部因子的条件。

\subsection{可去奇点}
\label{b3:subsec:030201}

\Needspace{6\baselineskip}
\begin{definition}\label{b3:def:03-removable}
设 \(a\in\mathbb C\)、\(R>0\)，\(f:D(a,R)\setminus\{a\}\to\mathbb C\) 全纯。若存在全纯函数 \(\widetilde f:D(a,R)\to\mathbb C\)，满足 \(\widetilde f(z)=f(z)\) 对所有 \(z\ne a\) 成立，则称 \(a\) 为 \(f\) 的\term[kequ-qidian]{可去奇点}[removable singularity]。
\end{definition}

延拓只能在中心补入一个值，连续性迫使这个值等于 \(\lim_{z\to a}f(z)\)，所以若延拓存在，它必唯一。真正的判别不必事先知道这个极限，只须给出整个去心邻域的上界。

\Needspace{6\baselineskip}
\begin{theorem}[可去奇点判别]\label{b3:thm:03-removable}
设 \(a\in\mathbb C\)、\(R>0\)，\(f:D(a,R)\setminus\{a\}\to\mathbb C\) 全纯。则以下条件等价：\(a\) 是可去奇点；以 \(a\) 为中心的 Laurent 展开没有非零负幂系数；存在 \(0<r<R\)、\(M<\infty\)，使 \(|f(z)|\le M\) 对 \(0<|z-a|<r\) 成立；有限极限 \(\lim_{z\to a}f(z)\) 存在。成立时，延拓在 \(a\) 处的值为 Laurent 系数 \(c_0\)。
\end{theorem}
\begin{proof}
若 \(f\) 可去，延拓函数在 \(D(a,R)\) 内的 Taylor 展开，在去心圆盘上也是原函数的 Laurent 展开。唯一性说明所有负幂系数为零。反之，若负幂系数全为零，Laurent 展开只有正幂部分，它在整个 \(D(a,R)\) 全纯，中心值为 \(c_0\)，直接给出延拓。

若存在题设的局部上界，则对每个 \(m\ge1\) 及 \(0<\rho<r\)，\cref{b3:prop:03-laurent-coefficients}给出
\[
 |c_{-m}|\le M\rho^m.
\]
系数与半径无关，令 \(\rho\downarrow0\) 得 \(c_{-m}=0\)。这对所有 \(m\) 成立，因此有界性推出可去。可去时，延拓在一个较小闭圆盘上连续，故局部有界。最后，可去蕴含有限极限存在；有限极限又使函数在充分小的去心圆盘内有界。四个条件由此等价。
\end{proof}

例如在零点以外考虑
\[
 f(z)=\frac{e^z-1-z}{z^2}.
\]
指数级数给出 \(f(z)=1/2+z/6+z^2/24+\cdots\)，这个幂级数在零点也收敛，所以补入 \(f(0)=1/2\) 后全纯。若只沿某一条路径看极限，结论就不可靠：\(e^{-1/z^2}\) 沿实轴趋于零，但在 \(z=\mathrm{i}y\) 时等于 \(e^{1/y^2}\)，沿虚轴无界。可去判别中的上界必须同时覆盖所有充分接近中心的方向。

\ProseParagraphBreak

有时表达式本身的上界不易取得，却知道乘以 \(z-a\) 以后比常数还小。这也足以排除全部负幂。

\Needspace{6\baselineskip}
\begin{corollary}\label{b3:cor:03-removable-small-growth}
设 \(a\in\mathbb C\)、\(R>0\)，\(f:D(a,R)\setminus\{a\}\to\mathbb C\) 全纯。若
\[
 \lim_{z\to a}(z-a)f(z)=0,
\]
则 \(a\) 是 \(f\) 的可去奇点。
\end{corollary}
\begin{proof}
令 \(g(z)=(z-a)f(z)\)。有限极限使 \(g\) 可去，并使延拓在 \(a\) 处取值零。它的 Taylor 展开因而没有常数项，可写成 \(g(z)=(z-a)h(z)\)，其中 \(h\) 在整个 \(D(a,R)\) 全纯。在 \(z\ne a\) 时 \(h(z)=f(z)\)，所以 \(h\) 是所需延拓。
\end{proof}

这里极限必须等于零。函数 \(1/(z-a)\) 满足 \((z-a)f(z)=1\)，却不能补成连续函数。乘去一个分母后得到非零极限，描述的正是下一小节的一阶极点。

\subsection{极点}
\label{b3:subsec:030202}

极点可以看作全纯函数零点的倒数：去掉一个有限整数次的分母后，剩余因子在中心全纯且非零。这一非零条件确定了真正的阶数。

\Needspace{6\baselineskip}
\begin{definition}\label{b3:def:03-pole}
设 \(a\in\mathbb C\)、\(R>0\)，\(f:D(a,R)\setminus\{a\}\to\mathbb C\) 全纯。若存在整数 \(m\ge1\)、\(0<r<R\) 以及全纯函数 \(g:D(a,r)\to\mathbb C\)，满足 \(g(a)\ne0\) 且
\[
 f(z)=(z-a)^{-m}g(z)\quad(0<|z-a|<r),
\]
则称 \(a\) 为 \(f\) 的 \(m\) 阶\term[jidian]{极点}[pole]。
\end{definition}

\Needspace{6\baselineskip}
\begin{proposition}\label{b3:prop:03-pole-criterion}
设 \(a\in\mathbb C\)、\(R>0\)，\(f:D(a,R)\setminus\{a\}\to\mathbb C\) 全纯。则 \(a\) 是极点，当且仅当 \(|f(z)|\to\infty\)（\(z\to a\)）；也当且仅当 Laurent 展开只有有限个非零负幂系数且至少有一个。若最低非零指数为 \(-m\)，则极点阶数恰为 \(m\)，并有
\[
 \lim_{z\to a}(z-a)^m f(z)=c_{-m}\ne0.
\]
反之，若对某个正整数 \(m\)，上式左端存在且为非零有限值，则 \(a\) 是 \(m\) 阶极点。
\end{proposition}
\begin{proof}
若 \(f(z)=(z-a)^{-m}g(z)\)、\(g(a)\ne0\)，连续性允许缩小半径，使
\[
 \frac{|g(a)|}{2}\le|g(z)|\le\frac{3|g(a)|}{2}.
\]
因此
\[
 \frac{|g(a)|}{2|z-a|^m}\le|f(z)|
 \le\frac{3|g(a)|}{2|z-a|^m},
\]
从而 \(|f(z)|\to\infty\)。将 \(g\) 的 Taylor 展开除以 \((z-a)^m\)，最低非零指数是 \(-m\)，系数为 \(g(a)\)，其他负幂只可能有有限个，并得到所述极限。

反之，若 \(|f(z)|\to\infty\)，则在一个充分小的去心圆盘内有 \(|f(z)|>1\)，所以 \(f\) 无零点，\(h=1/f\) 全纯并趋于零。可去奇点判别使 \(h\) 延拓到中心，且 \(h(a)=0\)。它在去心圆盘上处处非零，所以不恒为零。由\cref{b3:prop:01-zero-factorization}，缩小圆盘后可写
\[
 h(z)=(z-a)^m H(z),\quad m\ge1,\quad H(a)\ne0.
\]
进一步缩小圆盘使 \(H\) 无零点，取倒数便得
\(f(z)=(z-a)^{-m}/H(z)\)，这是极点分解。

若只有有限个非零负幂系数，令 \(-m\) 为其中最低指数，则 \((z-a)^m f(z)\) 的展开只有非负幂，并在中心取非零值 \(c_{-m}\)，也给出极点分解。最后，若 \((z-a)^m f(z)\) 有非零有限极限，它由可去奇点判别延拓为中心非零的全纯函数，仍得到同一分解。最低非零指数的唯一性同时保证极点阶数唯一。
\end{proof}

这里 \(|f(z)|\to\infty\) 指对每个 \(M>0\)，都有 \(\delta>0\)，使 \(0<|z-a|<\delta\) 时 \(|f(z)|>M\)。仅仅“在任意去心邻域无界”弱得多，不能代替这一条件。本性奇点也可能在一列点上任意大，却同时在另一列点上保持有限。

\ProseParagraphBreak

商式的分母为零时，应先比较分子与分母的零点阶数。设 \(g,h\) 在 \(a\) 附近全纯且都不局部恒零，令它们在 \(a\) 的阶数分别为 \(p,q\ge0\)，其中中心非零时阶数取零。零点分解给出
\[
 g(z)=(z-a)^pG(z),\quad h(z)=(z-a)^qH(z),
 \quad G(a)H(a)\ne0.
\]
在缩小的邻域内 \(H\) 无零点，故
\[
 \frac{g(z)}{h(z)}=(z-a)^{p-q}\frac{G(z)}{H(z)}.
\]
当 \(p>q\) 时，商延拓后有 \(p-q\) 阶零点；当 \(p=q\) 时，商的可去值为 \(G(a)/H(a)\ne0\)；当 \(p<q\) 时，商有 \(q-p\) 阶极点。分母的形式阶数只有在分子不消失时才直接等于极点阶数。

\ProseParagraphBreak

例如 \(1-\cos z=z^2/2+O(z^4)\)、\(\sin z=z+O(z^3)\)，所以
\[
 \frac{1-\cos z}{z\sin z}
 =\frac{1/2+O(z^2)}{1+O(z^2)}
\]
在零点可去，补入值为 \(1/2\)。而
\[
 \frac{1-\cos z}{(\sin z)^3}
 =\frac1z\frac{1/2+O(z^2)}{1+O(z^2)}
\]
在零点只有一阶极点，尽管分母具有三阶零点。这里各个 \(O(z^2)\) 都来自完整邻域内的收敛 Taylor 级数，分母中的剩余因子在中心等于一，故可以在小圆盘内求全纯倒数。

\ProseParagraphBreak

在开集 \(U\subseteq\mathbb C\) 上，若函数除一组在 \(U\) 内没有聚点的极点外全纯，并在各极点赋值 \(\infty\)，就称它为\term[ya-chun-hanshu]{亚纯函数}[meromorphic function]。极点分解说明它局部是两个全纯函数的商，分母不局部恒零；反过来，上面的公共零点消去也说明，这样的商补入可去值后局部只有极点。例如 \(\cot z=\cos z/\sin z\) 的分母在 \(\pi\mathbb Z\) 为单零点，分子在那里不为零，故这些点都是一阶极点。

\ProseParagraphBreak

对不局部恒零的亚纯函数，可以把阶数统一约定为：零点阶数取正，极点阶数取负，非零全纯点的阶数取零。局部分解立即给出乘积阶数相加、商的阶数相减。求和时最低项可能抵消，例如
\[
 \frac1z+\left(-\frac1z+z\right)=z.
\]
两项各有极点，和却有零点。只有比较首个非零系数以后，才能判断这种抵消是否发生。

\subsection{本性奇点}
\label{b3:subsec:030203}

\Needspace{6\baselineskip}
\begin{definition}\label{b3:def:03-essential}
设 \(a\in\mathbb C\)、\(R>0\)，\(f:D(a,R)\setminus\{a\}\to\mathbb C\) 全纯。若 \(a\) 既不是可去奇点也不是极点，则称它为 \(f\) 的\term[benxing-qidian]{本性奇点}[essential singularity]。
\end{definition}

按前两项判别，本性奇点恰好对应无穷多个非零负幂系数。三类因此互斥且穷尽：主部或者全为零，或者只有有限个非零项，或者有无穷多个非零项。中心是否具有某条路径上的极限，不在这三种代数条件中。

\ProseParagraphBreak

例如
\[
 e^{1/z}=\sum_{n=0}^{\infty}\frac{z^{-n}}{n!}
 \quad(z\ne0)
\]
的负幂无限多，所以零是本性奇点。沿正实轴趋近零时函数趋于无穷，沿负实轴则趋于零。沿虚轴也没有单一极限：取
\[
 z_k=\frac1{2\pi\mathrm{i}k},
 \quad w_k=\frac1{(2k+1)\pi\mathrm{i}}
 \quad(k\ge1),
\]
两列都趋于零，却分别满足 \(e^{1/z_k}=1\)、\(e^{1/w_k}=-1\)。而指数函数没有零点，所以值零始终不被取得。一个本性奇点附近可以同时出现增长、衰减和振荡。

\ProseParagraphBreak

孤立奇点分类的前提是完整去心圆盘上已经有一个单值全纯函数。平方根在零附近的困难不满足这一前提。若存在全纯函数 \(g\) 定义在整个 \(0<|z|<R\) 上并满足 \(g(z)^2=z\)，则 \(|g(z)|=|z|^{1/2}\)，可去奇点判别使它延拓到零，且 \(g(0)=0\)。延拓不恒为零，故其零点阶数是某个正整数 \(m\)，那么 \(g^2\) 的零点阶数为 \(2m\)，不可能等于 \(z\) 的阶数一。这证明全去心圆盘上不存在这样的分支；在割去一条射线的区域选取平方根，不与此矛盾。

\subsection{Casorati–Weierstrass 定理}
\label{b3:subsec:030204}

\term[casorati-weierstrass-dingli]{Casorati--Weierstrass 定理}[Casorati--Weierstrass theorem]说明本性奇点附近的像不能避开任何开圆盘。证明的构造来自一个反问：如果函数始终离某个值有正距离，那么差的倒数就有界，可以用可去奇点判别处理。

\Needspace{6\baselineskip}
\begin{theorem}[Casorati--Weierstrass]\label{b3:thm:03-casorati-weierstrass}
设 \(a\in\mathbb C\)、\(R>0\)，\(f:D(a,R)\setminus\{a\}\to\mathbb C\) 全纯，且 \(a\) 是本性奇点。则对每个 \(0<r<R\)，集合 \(f(D(a,r)\setminus\{a\})\) 在 \(\mathbb C\) 中稠密。等价地，对任意 \(b\in\mathbb C\)、\(\varepsilon>0\) 和 \(0<r<R\)，都存在 \(z\) 满足
\[
 0<|z-a|<r,\quad |f(z)-b|<\varepsilon.
\]
\end{theorem}
\begin{proof}
若对某个 \(0<r<R\) 像集不稠密，其闭包的补集含有一个开圆盘。因此存在 \(b\in\mathbb C\)、\(\varepsilon>0\)，使
\[
 |f(z)-b|\ge\varepsilon\quad(0<|z-a|<r).
\]
令 \(g(z)=1/(f(z)-b)\)。分母无零点使 \(g\) 全纯，并且 \(|g(z)|\le1/\varepsilon\)，所以 \(g\) 可以全纯延拓到 \(a\)。

若延拓满足 \(g(a)\ne0\)，连续性保证在一个小圆盘内 \(g\) 始终非零，于是 \(b+1/g\) 在整个小圆盘全纯，并在去心部分等于 \(f\)。这使 \(a\) 可去，与假设矛盾。

若 \(g(a)=0\)，则 \(g\) 在去心圆盘上非零，因而不局部恒零。零点分解给出
\[
 g(z)=(z-a)^m h(z),\quad m\ge1,\quad h(a)\ne0.
\]
缩小圆盘使 \(h\) 无零点，再写
\[
 f(z)=b+\frac1{(z-a)^m h(z)}
 =(z-a)^{-m}\left(b(z-a)^m+\frac1{h(z)}\right).
\]
括号内是全纯函数，在中心取非零值 \(1/h(a)\)，所以 \(a\) 是 \(m\) 阶极点，也与假设矛盾。两个可能的延拓值都被排除，故像集必须稠密。
\end{proof}

定理同时允许缩小自变量邻域和目标误差。对任意给定 \(b\in\mathbb C\)，取 \(r_k=\min\{R/2,1/k\}\)、\(\varepsilon_k=1/k\)，即可选择 \(z_k\ne a\)，使 \(z_k\to a\) 且 \(f(z_k)\to b\)。它并不保证某一点真正满足 \(f(z)=b\)，例如 \(e^{1/z}\) 就遗漏零值。\chapref{b3:ch:14}的 Picard 大定理将进一步讨论精确取值，而这里的证明只给稠密性。

\ProseParagraphBreak

任意去心邻域的像既然稠密，其中的模就没有有限上界，模的下确界又必为零。这两个结论都针对整个邻域，不能强加到每一条趋近中心的路径上。稠密性给出的离散点列，也不能自动连成一条使函数值收敛的连续路径：连接两点时，函数的像可能离开指定小圆盘。点列聚值、沿路径的极限和精确取值是不同的问题。

\ProseParagraphBreak

无穷远处的孤立奇点通过换坐标归入上述分类。设 \(R_0>0\)，\(f\) 在 \(|z|>R_0\) 全纯，令
\[
 w=\frac1z,\quad F(w)=f(1/w)
 \quad(0<|w|<1/R_0).
\]
约定 \(f\) 在无穷远可去、有 \(m\) 阶极点或本性，分别是指 \(F\) 在零点具有对应类型。这里分类的是函数值；本章下一节在变换积分微分时，还会出现换元导数，不能把两种计算混为一谈。

\Needspace{6\baselineskip}
\begin{proposition}[整函数在无穷远的分类]\label{b3:prop:03-entire-at-infinity}
设 \(f:\mathbb C\to\mathbb C\) 为整函数，通过 \(F(w)=f(1/w)\) 定义它在无穷远处的奇点类型。则无穷远可去当且仅当 \(f\) 为常函数；无穷远为 \(m\) 阶极点当且仅当 \(f\) 是次数恰为 \(m\ge1\) 的多项式；无穷远为本性奇点当且仅当 \(f\) 不是多项式。
\end{proposition}
\begin{proof}
若 \(F\) 在零点可去，它在小圆盘内有界，因此 \(f\) 在某个大圆之外有界。\(f\) 在该大圆所围闭圆盘上连续，也有界，合并两项上界便得到全平面有界。由\cref{b3:thm:02-liouville}，\(f\) 为常函数。反方向直接换坐标即可。

若 \(F\) 在零点有 \(m\) 阶极点，局部分解给出小 \(|w|\) 时 \(|F(w)|\le C|w|^{-m}\)。所以大 \(|z|\) 时有 \(|f(z)|\le C|z|^m\)；再利用紧圆盘上的有界性，增加常数后得到全平面的估计
\[
 |f(z)|\le C'(1+|z|)^m.
\]
\Cref{b3:prop:02-polynomial-growth}说明 \(f\) 是次数不超过 \(m\) 的多项式。它不能是常函数，否则 \(F\) 在零点可去。若其次数为 \(d\ge1\)，最高次系数为 \(a_d\ne0\)，则
\[
 F(w)=w^{-d}\bigl(a_d+a_{d-1}w+\cdots+a_0w^d\bigr).
\]
它在零点的极点阶数恰为 \(d\)，故 \(d=m\)。这个式子也直接证明任意 \(m\) 次多项式在无穷远有 \(m\) 阶极点。最后，由孤立奇点三类的穷尽性，余下的非多项式整函数恰对应本性奇点。
\end{proof}

指数函数和正弦函数的 Taylor 展开都有无穷多个非零系数，因此它们在无穷远都是本性奇点；\(z^m\) 在无穷远则有 \(m\) 阶极点。另一方面，\(1/z\) 换成 \(w\) 后为 \(w\)，在无穷远可去且取值零，却不是常函数，因为它在有限点零处不全纯。命题中的整性负责排除这种有限奇点，不能只由大圆外的衰减推出函数恒零。

% !TeX root = ../../Book3.tex
\section{留数}
\label{b3:sec:0303}

Laurent 展开的所有系数记录局部行为，闭曲线积分却只挑出 \(-1\) 次项。将多个孤立奇点围起来时，各个局部贡献相加，这使许多实积分也可以转化为有限次代数计算。

\subsection{留数}
\label{b3:subsec:030301}

\begin{definition}\label{b3:def:03-residue}
设 \(a\) 是全纯函数 \(f\) 的孤立奇点。其 Laurent 系数 \(c_{-1}\) 称为 \(f\) 在 \(a\) 的\term[liushu]{留数}[residue]，记为 \(\operatorname{Res}(f,a)\)。
\end{definition}
\symidx{\(\operatorname{Res}(f,a)\)：孤立奇点处的留数}
\begin{proposition}\label{b3:prop:03-residue-formulas}
若 \(a\) 是 \(m\) 阶极点，则
\[
 \operatorname{Res}(f,a)=\frac1{(m-1)!}\left.\frac{\mathrm{d}^{m-1}}{\mathrm{d}z^{m-1}}\bigl((z-a)^mf(z)\bigr)\right|_{z=a}.
\]
特别地，若 \(f=g/h\)，\(g,h\) 在 \(a\) 附近全纯，\(h(a)=0\)、\(h'(a)\ne0\)，则 \(\operatorname{Res}(f,a)=g(a)/h'(a)\)。
\end{proposition}
\begin{proof}
令 \(G=(z-a)^mf\)。其 Taylor 展开的 \((m-1)\) 次系数，除以 \((z-a)^m\) 后就是 \(-1\) 次系数。单零分母写成 \(h(z)=(z-a)H(z)\)，\(H(a)=h'(a)\)，故 \((z-a)f\) 在中心值为 \(g(a)/h'(a)\)。即使 \(g(a)=0\)，公式仍给留数0，此时奇点可能已经可去。
\end{proof}
对本性奇点仍可用级数提取留数，例如 \(e^{1/z}\) 在0的留数为1；有限阶求导公式却没有可用的整数 \(m\)。留数为零不意味着奇点可去，\(1/z^2\) 就是一例。

\ProseParagraphBreak

留数在合法的局部换元下属于同一个微分量。
\begin{proposition}[留数的坐标不变性]\label{b3:prop:03-residue-coordinate}
设 \(z=\psi(w)\) 是 \(b\) 附近的双全纯坐标，\(\psi(b)=a\)，而 \(f\) 在 \(a\) 有孤立奇点。则
\[
 \operatorname{Res}(f,a)=\operatorname{Res}\bigl(f(\psi(w))\cdot\psi'(w),b\bigr).
\]
\end{proposition}
\begin{proof}
取足够小的正向 \(w\) 圆，其像是包围 \(a\) 的正向简单闭曲线；全纯非零导数保持方向。曲线换元使两个微分积分相同。像曲线与一个小 \(z\) 圆在去心邻域之间可变形，Cauchy 定理使积分相同。两边各除以 \(2\pi\mathrm{i}\)，得到留数相等。
\end{proof}
因此只把 \(f(\psi(w))\) 拿来算留数通常是错误的，必须乘 \(\psi'(w)\)。例如 \(z=2w\)、\(f(z)=1/z\)，合成函数留数是 \(1/2\)，合成微分乘2后才是1。无穷远留数的 \(-w^{-2}\) 因子就是这一坐标规则的具体实例。

\ProseParagraphBreak

一个全纯函数的导数在任意孤立奇点处留数为零。将 Laurent 展开逐项微分，要产生 \(z^{-1}\) 项只能来自原来的 \(z^0\) 项，但它的导数系数为0。反向对同心环域也成立：留数为零时可逐项积分得到原函数。留数检测的是一个微分是否有非零局部周期，而不只是函数值增长快慢。

\subsection{留数定理}
\label{b3:subsec:030302}

\term[liushu-dingli]{留数定理}[residue theorem]先以简单围道版本建立。
\begin{theorem}[留数]\label{b3:thm:03-residue}
设 \(V\) 是有界 Jordan 区域，边界正向且分段光滑。若 \(f\) 在 \(\overline V\) 的邻域内除有限个内部孤立奇点 \(a_1,\ldots,a_N\) 外全纯，则
\[
 \int_{\partial V}f(z)\,\mathrm{d}z=2\pi\mathrm{i}\sum_{j=1}^N\operatorname{Res}(f,a_j).
\]
\end{theorem}
\begin{proof}
在每个奇点周围选互不相交、闭包包含于 \(V\) 的小圆盘。将挖去这些圆盘的区域沿有限条连接弧切开，再按有限小片细分；各小片上有局部原函数，其边界积分为零。内部边和切口两侧相消，留下外边界及顺时针小圆。故外边界积分等于各个逆时针小圆积分之和。小圆上 Laurent 展开一致收敛，逐项积分只有 \(-1\) 次项留下 \(2\pi\mathrm{i}\)，即得结论。若连接弧或边界不是多边形，可先在局部原函数圆盘中更换为多边形代表，端点公式保持积分。
\end{proof}
定理的有限性来自奇点与边界分离。若奇点向边界积聚，不能未经收敛核对就把有限和改成无穷和。自交闭曲线则需要把每个奇点的绕数作为权重，精确形式在\chapref{b3:ch:10}建立。

\subsection{实积分}
\label{b3:subsec:030303}

计算实积分时，先说明所选围道和附加弧积分为何消失。以
\[
 I=\int_{-\infty}^{\infty}\frac{\mathrm{d}x}{x^2+1}
\]
为例，在上半平面取半径 \(R>2\) 的半圆，内部仅有极点 \(\mathrm{i}\)，留数为 \(1/(2\mathrm{i})\)。半圆上的函数模不超过 \(1/(R^2-1)\)，弧长为 \(\pi R\)，其积分趋于0。实轴截断积分单调趋于反常积分，故 \(I=\pi\)。

\ProseParagraphBreak

若再有振荡因子 \(e^{\mathrm{i}tz}\)，\(t>0\)，在上半平面其模为 \(e^{-t\operatorname{Im}z}\le1\)。同一弧估计仍适用，留数变为 \(e^{-t}/(2\mathrm{i})\)，得到
\[
 \int_{-\infty}^{\infty}\frac{e^{\mathrm{i}tx}}{x^2+1}\,\mathrm{d}x=\pi e^{-t}.
\]
负 \(t\) 时选下半平面，或对称换元，得到 \(\pi e^{-|t|}\)。如果分母只衰减如 \(1/z\)，上述弧长估计并不趋于0，须另外证明振荡提供的衰减，不能直接援引“半圆消失”。

\ProseParagraphBreak

对 \(\int_0^{2\pi}R(\cos\theta,\sin\theta)\,\mathrm{d}\theta\)，置 \(z=e^{\mathrm{i}\theta}\)，则 \(\cos\theta=(z+z^{-1})/2\)、\(\sin\theta=(z-z^{-1})/(2\mathrm{i})\)、\(\mathrm{d}\theta=\mathrm{d}z/(\mathrm{i}z)\)。最后的微分因子不可遗漏，单位圆上有极点时则不能直接使用通常留数定理。

围道选择须同时控制奇点和大圆弧。对 \(\int_{\mathbb R}e^{\mathrm{i}ax}/(x^2+1)\,\mathrm{d}x\)，当 \(a>0\) 时闭向上半平面，因为 \(|e^{\mathrm{i}az}|=e^{-a\operatorname{Im}z}\le1\)；当 \(a<0\) 时应闭向下半平面，并记下顺时针的负号。若不随 \(a\) 改变半平面，大圆弧上的指数会增长，留数计算虽形式正确却不能代表原积分。最终两种选择都给 \(\pi e^{-|a|}\)，而 \(a=0\) 可以由有理函数估计直接处理。

\ProseParagraphBreak
实轴上有极点时还须说明积分的意义。通常的广义积分、Cauchy 主值以及绕开奇点的小半圆贡献并不相同。留数定理只计算闭围道积分，不会替读者选择其中一种意义。先标明每段截断积分和小圆弧的方向，再取极限，才能看清哪些项趋于零、哪些项留下半留数贡献。

\subsection{无穷远点}
\label{b3:subsec:030304}

无穷远点用坐标 \(w=1/z\) 检查函数奇点：研究 \(f(1/w)\) 在 \(w=0\) 的分类。留数属于微分 \(f(z)\,\mathrm{d}z\)，换元又给出额外的 \(-w^{-2}\)：
\begin{definition}\label{b3:def:03-infinity-residue}
若 \(f\) 在某个大圆之外全纯，定义
\[
 \operatorname{Res}(f,\infty)\coloneqq-\operatorname{Res}\left(\frac{f(1/w)}{w^2},0\right).
\]
\end{definition}
例如 \(f(z)=1/z\) 在无穷远作为函数可去，中心值为0，但微分的无穷远留数为 \(-1\)。函数的奇点分类和微分的留数不是同一问题。
\begin{proposition}\label{b3:prop:03-global-residue}
对有理函数 \(f\)，所有有限极点留数与无穷远留数之和为零。
\end{proposition}
\begin{proof}
选大圆包围全部有限极点，有限留数和为大圆积分除以 \(2\pi\mathrm{i}\)。换元 \(w=1/z\) 把逆时针大圆变成顺时针小圆，微分为 \(-w^{-2}\,\mathrm{d}w\)。根据定义，该积分是 \(-2\pi\mathrm{i}\operatorname{Res}(f,\infty)\)，故留数总和为零。
\end{proof}
计算无穷远留数时也可直接取 \(f\) 在大圆外 Laurent 展开的 \(z^{-1}\) 系数再加负号。例如 \(z^2/(z^3-1)\) 的这一系数为1，所以无穷远留数为 \(-1\)。

\ProseParagraphBreak

扩充平面上的亚纯函数实际上全部是有理函数。
\begin{theorem}\label{b3:thm:03-sphere-rational}
若 \(f\) 在扩充复平面上亚纯，则存在多项式 \(P,Q\)，\(Q\) 不恒零，使 \(f=P/Q\)。
\end{theorem}
\begin{proof}
极点在球面上离散，而球面紧，故极点有限；否则无限点列有聚点，与该处亚纯局部模型矛盾。先在每个有限极点减去其有限 Laurent 主部，所得函数 \(g\) 在整个有限平面全纯。在无穷远，原函数与减去的有理主部仍亚纯，故 \(g(1/w)\) 在0可去或有有限阶极点。于是存在整数 \(m\ge0\)，使大 \(|z|\) 时 \(|g(z)|\le C|z|^m\)。由第二章的增长版 Liouville 推论，\(g\) 是多项式。加回有限个有理主部，得到有理函数。
\end{proof}
这一定理利用了无穷远也亚纯的条件。整函数 \(e^z\) 在有限平面无极点，却在无穷远有本性奇点，所以不是反例。它也说明有理函数在球面只需有限多局部数据，整函数的无穷远增长则可能包含无限信息。

% !TeX root = ../../Book3.tex
\section{零点计数}
\label{b3:sec:0304}

函数值绕原点转了多少圈，可以通过 \(f'/f\) 的积分读出。零点和极点分别使辐角增加和减少；小扰动若不能把边界像推过原点，内部零点总数就不会变化。

\subsection{辐角原理}
\label{b3:subsec:030401}

对不经过原点的闭连续路径 \(\eta\)，可沿参数区间选连续辐角 \(\theta\)。具体地，把区间细分，使每段像位于一个不包含0的对数圆盘，在相邻端点给辐角加适当的 \(2\pi\) 整数以接合。因首尾函数值相同，\(\theta(b)-\theta(a)\) 必为 \(2\pi\) 的整数倍，其商称为\term[raoshu]{绕数}[winding number]。若路径分段光滑，局部对数求导后相加，得到
\[
 \operatorname{Ind}(\eta,0)=\frac1{2\pi\mathrm{i}}\int_\eta\frac{\mathrm{d}w}{w}.
\]
\term[fujiao-yuanli]{辐角原理}[argument principle]把该整数转成零点计数。
\begin{theorem}[辐角原理]\label{b3:thm:03-argument-principle}
设 \(f\) 在有界分段光滑 Jordan 区域 \(V\) 的闭包邻域亚纯，边界无零点也无极点。则
\[
 \frac1{2\pi\mathrm{i}}\int_{\partial V}\frac{f'(z)}{f(z)}\,\mathrm{d}z=N-P,
\]
其中 \(N\) 与 \(P\) 分别是内部零点和极点的总重数。
\end{theorem}
\begin{proof}
紧性与边界无零极点使内部只有有限个零极点。在零点 \(a\) 的 \(m\) 阶分解中，\(f=(z-a)^mg\)，故 \(f'/f=m/(z-a)+g'/g\)，留数为 \(m\)。在 \(m\) 阶极点处同样得到留数 \(-m\)。留数定理相加给出公式；链式法则说明左端也等于像曲线 \(f\circ\partial V\) 绕0的绕数。
\end{proof}

\subsection{Rouché 定理}
\label{b3:subsec:030402}

\term[rouche-dingli]{Rouché 定理}[Rouché's theorem]允许把难计数的函数与简单函数比较。
\begin{theorem}[Rouché]\label{b3:thm:03-rouche}
设 \(f,g\) 在有界分段光滑 Jordan 区域 \(V\) 的闭包邻域全纯。若边界处 \(|g|<|f|\)，则 \(f\) 与 \(f+g\) 在 \(V\) 内的零点总重数相同。
\end{theorem}
\begin{proof}
对 \(0\le t\le1\)，令 \(f_t=f+tg\)。边界上 \(|f_t|\ge|f|-|g|>0\)，紧边界使右端有共同正下界。辐角原理给出的数
\(N(t)=(2\pi\mathrm{i})^{-1}\int_{\partial V}f_t'/f_t\,\mathrm{d}z\)
依 \(t\) 连续，又总为整数，故在区间上为常数。取 \(t=0,1\) 即得结论。
\end{proof}
例如在 \(|z|=2\) 上，\(|z^5|=32>2+1\)，故 \(z^5+z+1\) 在盘内有5个零点。严格不等式排除了整个变形中边界出现零点。若只有 \(|g|\le|f|\)，边界可能产生零点，开域中的计数就可能改变。

\ProseParagraphBreak

Rouché 的比较方向由边界决定，不能只看多项式的次数。在 \(|z|=1\) 上，\(z^4+3z+1\) 应比较 \(3z\)，因为余项模不超过2；在 \(|z|=2\) 上，\(z^4\) 的模为16，大于 \(3|z|+1=7\)，应比较最高次项。因此同一个多项式在单位盘有1个零点，在半径2盘有4个零点，其余3个位于两圆之间，均计重数。不同半径的零点计数可以相减定位环域中的零点，但边界必须各自无零点。

\ProseParagraphBreak

参数变化时，若一个紧参数区间内的边界函数始终无零点，零点数就保持不变。共同正下界来自参数区间与边界乘积的紧性。如果参数范围不紧，最低模可能趋于0，即使每个固定参数都无边界零点，也不能得到统一的扰动余量。实际使用时可局部按参数取一个小区间，再连通地传播整数计数；不必先找到一个覆盖所有无界参数的误差常数。

\subsection{重数}
\label{b3:subsec:030403}

零点个数若不计重数，就不具备稳定性。函数 \(z^2-\varepsilon\) 在 \(\varepsilon\ne0\) 时有两个不同零点，极限 \(z^2\) 只有一个零点，然而总重数始终为2。
\begin{proposition}[零点局部稳定]\label{b3:prop:03-local-zero-stability}
设 \(f\) 在 \(a\) 处有 \(m\) 阶零点。在充分小的闭圆盘 \(\overline{D(a,r)}\) 上，若全纯函数 \(h\) 满足边界估计 \(|h-f|<\min_{|z-a|=r}|f(z)|\)，则 \(h\) 在圆盘内恰有 \(m\) 个零点，计重数。
\end{proposition}
\begin{proof}
由孤立零点定理选圆盘，使 \(f\) 在其中只有中心零点且边界无零点。给定边界估计就是 Rouché 定理对 \(f\) 与 \(h-f\) 的条件，故两者零点总重数相同。
\end{proof}
若希望定位单个零点，先在每个已知零点周围选互不相交的小圆盘，再分别估计。只在一个大圆上计总数，不会自动给出各个零点的精确位置。

Rouché 定理中的严格不等式只需在边界成立，不要求内部成立。设 \(p(z)=z^5+10z+1\)，在单位圆上以 \(10z\) 为主项，因为 \(|z^5+1|\le2<10\)，可知内部有一个零点；在半径2的圆上，\(|z^5|=32>21\ge|10z+1|\)，内部有五个零点。因此其余四个零点落在 \(1<|z|<2\)。这比求根公式更直接，并且适用于无法写出根的高次多项式。

\ProseParagraphBreak
若把严格不等式放宽为处处等号允许，结论会失败。单位圆上 \(|z|=|1|\)，但 \(z\) 在内部有一个零点，常数1没有零点。严格性保障整个变形 \(f+t g\) 在边界无零；一旦某个参数的零点经过边界，整数零点数便可能跳变。实际计算通常先在紧边界上寻找一个正的间隙，再允许小参数扰动。这个间隙是解析问题中控制根的稳定性的数值。

\subsection{对数导数}
\label{b3:subsec:030404}

对数导数还可以提取带权的零点信息。
\begin{proposition}\label{b3:prop:03-weighted-zeros}
在辐角原理的条件下，若 \(h\) 在 \(\overline V\) 的邻域全纯，则
\[
 \frac1{2\pi\mathrm{i}}\int_{\partial V}h(z)\cdot\frac{f'(z)}{f(z)}\,\mathrm{d}z
 =\sum_a m_a h(a)-\sum_b n_b h(b),
\]
两和分别遍历内部零点与极点，\(m_a,n_b\) 为阶数。
\end{proposition}
\begin{proof}
对数导数在 \(a\) 的主部只有 \(m_a/(z-a)\)，乘全纯 \(h\) 后留数为 \(m_ah(a)\)；极点处相同计算带负号。直接应用留数定理。
\end{proof}
取 \(h(z)=z\) 得到零点之和与极点之和的差；取 \(h(z)=z^k\) 得到幂和。若 \(f\) 在整个区域无零点，\(f'/f\) 本身全纯；它有没有全局原函数，还须看所有闭曲线周期是否为零。\chapref{b3:ch:05}将用这一条件构造全纯对数。

局部零点稳定性始终按重数计数。不同简单根可以合成重根，重根也可分裂；只有在分别选定无零边界的小圆盘之后，才有可持续追踪的局部零点总数。


\begin{chaptersummary}
环域的内外边界分别产生 Laurent 展开的负幂和正幂。可去奇点、极点与本性奇点由主部的三种形态区分，Casorati--Weierstrass 定理展示本性奇点的局部像。围道积分只保留各奇点的留数；对数导数又把零点和极点的阶数转换为留数。Rouché 定理依赖边界变形始终避开零点，因此稳定的是总重数而不是不同零点的个数。
\end{chaptersummary}

\BookExercises

% 本章题组自拟；以正文已建立的工具设计计算、反例、构造与综合训练。
\begin{exercise}\label{b3:ex:03-1}
分别在 \(0<|z|<1\)、\(1<|z|<2\)、\(|z|>2\) 展开 \(1/((z-1)(z-2))\)。
\end{exercise}
\begin{solution}
先分解为 \(-1/(z-1)+1/(z-2)\)。内环展开为 \(\sum_{n\ge0}(1-2^{-n-1})z^n\)；中环为 \(-\sum_{n\ge0}z^{-n-1}-\sum_{n\ge0}z^n/2^{n+1}\)；外环为 \(\sum_{n\ge0}(2^n-1)z^{-n-1}\)。各几何级数分别按所给半径条件收敛。
\end{solution}

\begin{exercise}\label{b3:ex:03-2}
分类零点处 \(\sin z/z\)、\((e^z-1)/z^3\)、\(e^{1/z}-1/z\) 的奇点并求留数。
\end{exercise}
\begin{solution}
第一项可去，留数0。第二项展开为 \(z^{-2}+\tfrac12z^{-1}+\tfrac16+\cdots\)，是二阶极点，留数 \(1/2\)。第三项负幂仍无限，但 \(z^{-1}\) 项被消掉，是本性奇点，留数0。
\end{solution}

\begin{exercise}\label{b3:ex:03-3}
若 \(f\) 在去心圆盘全纯并满足 \(|f(z)|\le C|z|^{-3/2}\)，证明它在0最多有一阶极点。
\end{exercise}
\begin{solution}
负指标系数满足 \(|c_{-m}|\le C\rho^{m-3/2}\)。对 \(m\ge2\)，令 \(\rho\downarrow0\) 得系数为零；因此主部只可能有 \(c_{-1}/z\)。该系数非零时是一阶极点，为零时可去。
\end{solution}

\begin{exercise}\label{b3:ex:03-4}
计算 \(\int_{-\infty}^{\infty}(x^2+1)^{-2}\,\mathrm{d}x\)。
\end{exercise}
\begin{solution}
在上半圆内只有 \(\mathrm{i}\) 的二阶极点。留数是 \(\left.(\mathrm{d}/\mathrm{d}z)(z+\mathrm{i})^{-2}\right|_{z=\mathrm{i}}=-2/(2\mathrm{i})^3=1/(4\mathrm{i})\)。弧估计不超过 \(\pi R/(R^2-1)^2\)，趋于0。因此积分为 \(2\pi\mathrm{i}/(4\mathrm{i})=\pi/2\)。
\end{solution}

\begin{exercise}\label{b3:ex:03-5}
证明 \(z^4+3z+1\) 在单位圆盘内恰有一个零点。
\end{exercise}
\begin{solution}
在单位圆上 \(|z^4+1|\le2<3=|3z|\)。用 \(3z\) 作为比较函数，Rouché 给出零点总重数为1；因此该零点也必为单零点。
\end{solution}

\begin{exercise}\label{b3:ex:03-6}
设 \(P\) 是次数 \(n\ge1\) 的多项式，所有零点都在 \(|z|<R\)。用围道积分表达这些零点之和与平方和，计重数。
\end{exercise}
\begin{solution}
分别取加权辐角原理中的 \(h(z)=z\) 与 \(z^2\)，得到 \((2\pi\mathrm{i})^{-1}\int_{|z|=R}zP'/P\,\mathrm{d}z\) 和 \((2\pi\mathrm{i})^{-1}\int_{|z|=R}z^2P'/P\,\mathrm{d}z\)。多项式无有限极点，所以无需减极点项。
\end{solution}

\begin{exercise}\label{b3:ex:03-7}
求 \(z^2/(z^3-1)\) 在三个有限极点以及无穷远的留数。
\end{exercise}
\begin{solution}
每个三次单位根 \(a\) 处，单零分母公式给出 \(a^2/(3a^2)=1/3\)。大圆外展开首项为 \(1/z\)，故无穷远留数为 \(-1\)。总和为 \(3\cdot1/3-1=0\)。
\end{solution}

\begin{exercise}\label{b3:ex:03-8}
设 \(f\) 在 \(0<|z|<1\) 全纯且有本性奇点。证明存在趋于0的点列 \(z_n\)，使 \(f(z_n)\rightarrow b\)，其中 \(b\) 是任意给定复数。是否能由此保证 \(f(z_n)=b\)？
\end{exercise}
\begin{solution}
对每个 \(n\)，Casorati--Weierstrass 稠密性在 \(0<|z|<1/n\) 中给出 \(|f(z_n)-b|<1/n\)。所以同时有 \(z_n\rightarrow0\) 与函数值趋于 \(b\)。精确取值不能保证，\(e^{1/z}\) 永远不取0。
\end{solution}
